\section{Literature Review}

\subsection{Clinical Documentation Burden and AI-Assisted Transcription}
Administrative documentation is a leading driver of clinician burnout and reduced care quality. Manual note-taking during consultations is time-consuming, error-prone, and diverts clinician attention from the patient. Systems that combine automatic speech recognition (ASR) with speaker diarization offer a path toward automated, speaker-labeled clinical transcripts, significantly reducing the documentation burden \cite{klusty2024clinical, saadat2025enhancing}.

Systematic reviews of AI-assisted clinical documentation show that integrating NLP, ASR, and large language models (LLMs) reduces transcription errors, supports real-time note generation via SOAP templates, and improves EHR interoperability. Techniques such as ClinicalBERT and GPT-4 fine-tuning, alongside hybrid extractive-abstractive summarization, have been demonstrated to accelerate structured note generation. Key limitations remain, however: hallucination risk and low clinician trust continue to hinder adoption \cite{saadat2025enhancing}.

Open-source generative LLMs, including Llama-3.3-70B, Phi-4-14B, and Qwen-2.5-14B, have been evaluated in zero-shot settings for clinical information extraction from medical texts. These models demonstrate strong performance, matching or exceeding fine-tuned RoBERTa baselines on structured inference tasks. Notably, native-language processing must be preserved; translation to English consistently degrades extraction quality, underscoring the importance of multilingual model support \cite{builtjes2025leveraging}.

\subsection{Automatic Speech Recognition for Clinical Speech}
Radford et al.\ introduced Whisper, a general-purpose ASR system trained on 680,000 hours of weakly supervised multilingual audio. Its zero-shot performance approaches human-level accuracy across diverse acoustic conditions, making it a compelling foundation for clinical pipelines that cannot rely on large labeled datasets \cite{radford2023whisper}.

Independent evaluation of general-purpose and medical-specialized ASR systems on 36 patient-clinician conversations reports word error rates (WER) of 8.8–10.5\% and word-level diarization error rates (WDER) of 1.8–13.9\% under controlled recording conditions. Counterintuitively, general-purpose models sometimes outperform medical-specialized variants on diarization accuracy. Medical concept recall (F-recall \textasciitilde{}0.48–0.49) remains low despite high transcription precision, motivating further LLM post-processing stages \cite{tran2022asr}.

To address domain-specific terminology gaps, United-MedASR layers a BART-based semantic correction model over Faster Whisper, fine-tuned on synthetic vocabulary derived from ICD-10, MIMS, and FDA databases. This approach achieves WER below 1\% on benchmark corpora without requiring private patient data, demonstrating that synthetic domain adaptation is a viable and privacy-preserving strategy \cite{banerjee2024medasr}.

PriMock57 provides 57 mock primary care consultations with manual utterance-level transcriptions and consultation notes, representing one of the few publicly available clinical audio benchmarks due to stringent privacy constraints. It enables reproducible evaluation of ASR and note-generation systems and serves as the evaluation corpus for this study \cite{papadopoulos2022primock57}.

Diarization accuracy is sensitive to audio quality, overlapping speech, and the choice of ASR tool. Fine-tuned LLMs applied as a post-processing layer can markedly improve diarization accuracy; however, models fine-tuned on one ASR system's outputs underperform when applied to others, motivating ensemble approaches. LLM post-processing with chain-of-thought (CoT) prompting further improves medical-concept WER and speaker role assignment in primary care transcripts \cite{efstathiadis2024diarization, adedeji2024sound}.

\subsection{Clinical Information Extraction and NLP Pipelines}
Empirical comparison of structured and unstructured EHR data on 172 clinical phenotypes from MIMIC-III shows that Doc2Vec-encoded textual features outperform structured measurements alone in 84\% of cases, consistently outperforming Bag-of-Words and Bag-of-Concepts representations. Combining both modalities yields statistically significant AUC improvements for 51 phenotypes, particularly where radiology and surgical reports carry decisive diagnostic context; for one exception (diabetes without complications), structured data alone performs better. These findings establish a clear empirical basis for multimodal pipelines that fuse structured and unstructured clinical data \cite{moldwin2021phenotyping}.

Hong et al.\ introduced a FHIR-based type system for integrating structured and unstructured EHR data via UIMA-based NLP tools (MedXN, MedTime), achieving F-scores of 0.73--0.99 across FHIR element representations on Mayo Clinic medication data \cite{hong2017integrating}. Building on this, the NLP2FHIR pipeline extended FHIR-based normalization into a scalable framework with 30 mapping rules, 62 normalization rules, and 11 FHIR extensions, covering Condition, Procedure, MedicationStatement, and FamilyMemberHistory resources with F-scores of 0.69--0.99 \cite{hong2020nlp2fhir}.

More recently, transformer-based NER combined with ontology-grounded concept normalization (SNOMED-CT, ICD-10, RxNorm, LOINC) and relation extraction has been shown to produce FHIR-compliant patient digital twin representations from free-text clinical narratives. Evaluation on MIMIC-IV Demo achieves 0.89 NER F1, 0.81 relation extraction F1, 91\% semantic completeness, and an interoperability score of 0.88. Concept normalization is identified as the single factor with the highest isolated impact on overall interoperability \cite{brens2025semantic}.

\subsection{FHIR-Based Interoperability and LLM-Driven Resource Synthesis}
HL7 FHIR Release 4B defines a modular, resource-based data model adopted by major healthcare regulators globally, and its RESTful architecture has driven broad adoption across EHR vendors and research platforms \cite{hl7fhir2022r4b}. FHIR-GPT demonstrates that LLMs can convert clinical narratives into FHIR medication statements with greater than 90\% exact match rate, outperforming rule-based NLP pipelines by 3--50\% across attributes such as route, dose, reason, form, and timing, confirming that LLMs can replace modular NLP pipelines for FHIR transformation while remaining flexible to varied input formats \cite{li2024fhirgpt}.

The Inferno framework advances this direction by using LLM agents equipped with SNOMED CT tools and code execution to synthesize FHIR resources from discharge letters, competing with a human baseline. Programmatic FHIR object construction via the \texttt{fhir.resources} library, rather than free-form JSON generation, reduces hallucination rates to below 2.3\%, illustrating the value of structured output constraints in agentic pipelines \cite{frei2025inferno}.

The SMART Text2FHIR pipeline demonstrates an NLP-centric alternative, using Apache cTAKES to extract clinical concepts from EHR notes and map them to FHIR Observation, Condition, Procedure, and MedicationStatement resources, with extensions for negation, NLP sourcing, and provenance. Increased FHIR data liquidity driven by interoperability regulations makes such NLP-to-FHIR pipelines viable at institutional scale \cite{miller2023smart}.

Finally, retrieval-augmented generation (RAG)-enhanced LLMs, including GPT-4o and Llama 3.2 405B, applied to MIMIC-IV structured data achieve 67–73\% attribute-level FHIR mapping accuracy. Detailed JSON schemas and semantic clustering help narrow ambiguity, while iterative prompting further refines outputs. Persistent challenges include model hallucinations and incomplete source field descriptions, highlighting the continued need for validation layers in automated FHIR generation pipelines \cite{riquelme2025llm}.

\subsection{Critical Synthesis and Research Gap}

Reviewed works collectively advance the constituent technologies of clinical audio-to-FHIR pipelines, yet several recurring weaknesses motivate the present study. Existing work treats ASR, clinical NLP, and FHIR generation as isolated research problems rather than interdependent stages of a single system: ASR studies evaluate transcription quality without propagating errors to downstream structuring tasks \cite{tran2022asr, radford2023whisper, adedeji2024sound}; NLP-to-FHIR pipelines assume clean text input \cite{hong2020nlp2fhir, brens2025semantic, miller2023smart}; and FHIR synthesis systems presuppose structured or semi-structured input rather than conversational speech \cite{li2024fhirgpt, frei2025inferno, riquelme2025llm}. Text-to-FHIR studies are evaluated on pre-existing EHR notes (NLP2FHIR on Mayo Clinic records, Inferno on discharge letters, SMART Text2FHIR on cTAKES-processed notes), inputs that are already partially structured and free of the overlapping speakers, informal register, and ASR noise present in real clinical audio. No prior work measures the cumulative effect of chaining these stages end-to-end.

Evaluation scope is further constrained by narrow resource coverage and proprietary data. Most FHIR generation studies target a single resource type: FHIR-GPT focuses on medication statements \cite{li2024fhirgpt} and Inferno on discharge-letter entities \cite{frei2025inferno}. The 67--73\% attribute-level accuracy reported for RAG-enhanced LLMs \cite{riquelme2025llm} illustrates a substantial gap even on clean structured input, a figure that would likely worsen on upstream ASR output. Hallucination and conformance errors persist across reviewed systems; Inferno reports sub-2.3\% hallucination rates only after applying constrained programmatic FHIR construction \cite{frei2025inferno}, and no reviewed study performs full HL7 FHIR R4B conformance validation on pipeline output derived from spoken audio.

These gaps collectively motivate a unified pipeline that: (a) begins from raw audio rather than pre-existing text; (b) integrates ASR, diarisation, LLM post-processing, clinical NLP, and FHIR generation as a single evaluated system; (c) uses PriMock57, a public and realistic clinical audio benchmark \cite{papadopoulos2022primock57}, to ensure reproducibility; and (d) includes HL7 FHIR R4B conformance validation as a first-class evaluation criterion.
